\section{Tính kiểm chứng được của khoa học}

\

Hãy để ý về định nghĩa của giả thuyết. Một khẳng định để được coi là một giả thuyết, thì nó phải có thể được kiểm chứng. Nói cách khác, mệnh đề đó cần phải được phát biểu sao cho có một phương pháp nào đó có thể kiểm tra được tính đúng đắn và qua đó chấp nhận hoặc bác bỏ giả thuyết.

Ngoài ra, kèm theo việc có thể xác định được tính đúng sai, nó còn cần phải có khả năng dự đoán các sự vật, hiện tượng chưa quan sát được. Trong quá khứ của thiên văn học, hai nhà toán học khi nghiên cứu chuyển động của sao Thiên Vương nhận thấy rằng nó di chuyển không đúng so với quỹ đạo được dự đoán. Sử dụng định luật vạn vật hấp dẫn của Niu-tơn, họ đã suy ra rằng cần có một hành tinh khác để tác động hấp dẫn lên sao Thiên Vương và làm nó đi chệch quỹ đạo. Thật là như vậy, khi họ chĩa kính viễn vọng vòng vị trí được dự đoán của hành tinh mới đó, họ đã tìm ra sao Hải Vương. Sự kiện này càng kiểm chứng tính chính xác của định luật.

Tuy nhiên, một định luật không phải cứ đúng nhiều lần là sẽ tuyệt đối đúng. Được khích lệ bởi thành công trong việc khám phá sao Hải Vương, các nhà thiên văn học đã thử thủ thuật tương tự với sao Thủy. Quỹ đạo của nó cũng cho thấy những bất thường, vì vậy họ đã đưa ra giả thuyết về một hành tinh ẩn khác nằm bên trong quỹ đạo của sao Thủy. Song, dù đã trải qua rất nhiều lần tính toán và tìm kiếm nhưng họ không thể phát hiện ra hành tinh đó. Cuối cùng, giả thuyết này bị bác bỏ sau khi Anh-xtanh đưa ra thuyết tương đối rộng, giải thích được sự bất thường trong quỹ đạo của sao Thủy\footnote{Bạn đọc muốn tìm hiểu thêm có thể xem xét tài liệu\cite{baum1997search}.}. Điều này cũng chứng tỏ rằng định luật của Niu-tơn không phải lúc nào cũng chính xác, mà chỉ áp dụng được trong giới hạn của cơ học cổ điển. Mặc dù vậy, bạn đọc cần phải tránh việc phủ nhận hay xem nhẹ khả năng ứng dụng phát hiện của Niu-tơn. Trong các hành trình khám phá vũ trụ của con người, định luật vạn vật hấp dẫn vẫn được sử dụng.

Tóm lại, mọi mệnh đề khoa học đều phải có thể được kiểm chứng và có khả năng dự đoán, nhưng chúng ta phải thực hiện đánh giá các mệnh đề đó theo cách \defText{biện chứng}. Nói cách khác, cần phải nhìn nhận một phát biểu trong mối quan hệ với các sự vật và hiện tượng khác, trong một không gian và thời gian nhất định, trước khi chúng ta đưa ra đánh giá về chúng.